\chapter{制度中的辩证法}

黑格尔意义上的法哲学既不是一种实定法学理论，也不是阐述一系列个体自由权利的自然法的原则学说。一般说来，黑格尔法哲学的对象并不是任何历史的或者超历史的权利(Rechten)体系，而就是（单数的）法(Recht)——更准确地说：法的“概念”及其“定在”或者“现实化”（第1节）。黑格尔谈到了“法的理念”，因为正如这部著作开篇所言，哲学探究的是理念，而不是人们通常所言的“单纯概念”。

\begin{quote}
相反地，哲学应该指出单纯概念的片面性和非真理性，同时指出，只有概念才具有现实性，并从而使自己现实化。\footnote{《法哲学原理》，第1节，附释。《全集》第7卷，H. 格洛克纳出版（1927年之后），第38页。（中译文引自〔德〕黑格尔：《法哲学原理》，范扬、张企泰译，商务印书馆1961年版，第17页。——译者）}
\end{quote}

这种言说方式以对一些“逻辑学”的基本术语规定的理解为前提。按照黑格尔的理论，“定在”意味着“规定了的存在”，在其“形态”和“实存”的规定性中，定在并不等同于日常语言所言的“现实”。黑格尔把前述“形态”和“实存”称作过渡为他者的直接事物的诸形式。历史上的实定法，比如说一个地区或国家的法律，是现存的和有效的。也就是说，它们“在历史上实际存在”。但是，《逻辑学》引入的“实存”这一术语乃是“存在和反思的直接统一”，因而界定的并不是现实的结构，而是现象的结构。因此，尽管法也必须显现出来，尽管“伦理事物的现象世界”在法哲学的概念序列中占据了突出的地位，然而这里现象的语言在方法上的作用是次要的。每一种现象都是直接的，正如同黑格尔本人所言，它来自根据，又回到根据\footnote{“回到根据”德文原文为“geht zu Grunde”，意为“死亡、毁灭”。——译者}。只有那种不再“过渡到他者”的根据才对逻辑的东西和历史的东西“进行中介”。在黑格尔的术语中，被中介者的形式所特有的这个根据，就叫“现实”：

\begin{quote}
现实事物是上述那种直接统一的设定存在，是达到了自身同一的关系；因此，它得免于过渡……它的定在只是它自己本身的表现，而非他物的表现。\footnote{《哲学科学百科全书》，第一部分，逻辑学，第142节，附释，《全集》第8卷，第320页。（中译文引自〔德〕黑格尔：《小逻辑》，贺麟译，商务印书馆1980年版，第295页。按德文原注标明此句引文出自“说明”，疑误。——译者）}
\end{quote}

法哲学的对象是法的现实，依照黑格尔的观点，法的现实在“自由的理念”之中，后者建立在对立规定的统一之中。历史上设定的法所“表现”的，是那种自己赋予自己“现实性”的“概念”，它仿照自然法，在“其为自由的意志”（第4节）的自由的规范原则下构成了“哲学法学”的出发点。哲学法学在方法上的创新并不在于从规范性的法原则推演出超历史的法规范体系，而是把通达理念的道路理解为各种制度的历史形成过程（这是一个辩证的——矛盾的历程），并且把这条曲折的道路理解为认识逻辑概念自身辩证法的必要条件。

\section*{1}

在法哲学中，“制度”这个词并不是一个规范的术语，而是从语言的实际使用拿来并在不同的语境中加以使用的。我们可以对其作狭义(a)和广义(b)之分。狭义的制度(a)借助罗马法的规定，而黑格尔只把罗马法解释为私法（就历史而言，这是错误的）。按照这个词的原初法学语境，“制度”的例子有作为罗马私法概念之基础的“罗马家父权、罗马婚姻状况”\footnote{《法哲学原理》，第3节，附释，《全集》第7卷，第43页。（中译文参见《法哲学原理》，第5页。——译者）}。这些例子的语境为黑格尔所言的（以财产和契约为基本制度的）“抽象法”领域，在《法哲学原理》概念构建的体系进展中，它超出了这一语境。对于黑格尔而言，只有（在“道德”一章中）批判了这些制度之法的历史局限性并在伦理的概念和定在中扬弃了这些制度之后，自由的理念才是“现实的”。伦理作为“活的善，在自我意识中具有它的知识和意志，通过自我意识的行动而达到它的现实性；另一方面自我意识在伦理性的存在中具有它的绝对基础和起推动作用的目的”（第142节）。[由此，伦理]也需要历史的形态。伦理的语言不外就是“由于自由的理念而是必然的、因此是现实的那些关系在它们全部范围内即在国家中的发展”\footnote{同上书，第148节，《全集》第7卷，第229页以下。（中译文引自《法哲学原理》，第167页。——译者）}。

在伦理及其与“国家法”——公法——的关系的语境下，制度概念获得了其更加广泛的意义(b)。黑格尔在这里当然不再继续使用拉丁化的表达，而是代之以德语同义词“Einrichtung”。伦理事物的“固定的内容”，那些摆脱了个体的主观偏好和意见的“伦理力量”（第145节），乃是“自在自为地存在着的法律和制度”\footnote{《法哲学原理》，第144节，《全集》第7卷，第227页。（中译文参见《法哲学原理》，第164页。——译者）}。

如前述，逻辑的东西和历史的东西在法中的中介导向“逻辑学”的现实概念，而现实概念也为自然法、实定法以及国家学之间的关系且由此而为现在即将进一步研究的制度的辩证法之起源奠定了基础。“现实”一词每个人都很熟悉，它的每一用法都是开放的，遗憾的是，对于这个具有如此多意义的词语必然引发的误解，黑格尔并没有进行防备。尽管《法哲学原理》“序言”在其关于现实的东西的合理性和合理性的东西的现实性的著名的或者说声名狼藉的格言中已经涉及“现实”一词的逻辑规范意义，但由于它又同时影射了一种前哲学的——通俗的语言运用，引起了对他的双重误解：一方面，他为实定的东西特别是历史中现成的（“普鲁士”）国家进行明确辩护，另一方面，他明确地拒斥自然法原则。确实，黑格尔对“哲学法”的构想与其近代自然法前辈的设想没有多少共同之处，但他与那种他的后辈和批判者都想要读出来的复辟时期的“国家哲学”计划则更是格格不入。

根据《逻辑学》的概念定义，“现实”意味着“本质与实存或内与外所直接形成的统一”\footnote{《哲学全书》，第142节，《全集》第8卷，第320页。（中译文引自《小逻辑》，第295页。按德文原注标明此句引文出自“说明”，疑误。——译者）}。《法哲学原理》中有两个环节与这种把现实的东西的结构规定为中介的用法相符。这两个环节属于法概念的辩证法，同时也属于它的定在、伦理制度中的现实化。我们采用稍微不同于黑格尔的术语，称之为法的反思性和实定性环节。“反思性”这一术语在这里表示的是与法的“基地”的关联，这种关联属于法的现实性。对法来说，这种关联是一个“精神物”——也就是法主体——的构成性的自身相关，这个主体同时知道自己是那些承认和遵守规范和法律的义务的承担者。黑格尔的“自由意志”概念是那种反思性的基本表达。按照作为“外化”与自身反思之统一的一切现实的辩证法，自由意志在它的他者中意求它自身，因此也就是意求“自由意志”（第27节）。“自由意志”，即人之为人的普遍的法权能力（Rechtsfähigkeit），就是近代自然法意义上的“法”的基本概念。但是，正如辩证的重构所表明的那样，此种意志的内容是制度的各种历史形态。《法哲学原理》并不认为有从一种规范的普遍的法权原则中演绎出法的那些特定内容的可能性。黑格尔把这种思想斥为一种抽象的自然法假设：“可以有一个法体系和一种应当是纯粹理性的——仅仅是理性的——法状态。”\footnote{《法哲学讲演录》第2卷，伊尔廷（K.-H. Ilting）出版，斯图加特-巴德·坎施塔特，1974年，第89页。}在制度的特定发展阶段——诸如罗马法或中世纪封建法——上的法的内容是偶然的历史形态的产物。尽管对这些制度不仅要从历史的角度进行描述和解释，而且也要按照人之为“自由人”的法权能力原则进行评价（第3节），《法哲学原理》还是非常一贯地禁止从自然法原则出发进行演绎。

理解“哲学的法”之辩证法的关键在于“实定性”概念。众所周知，“实定性”这一术语在黑格尔青年时期作品中居于中心地位，而在这里，它表达的是法概念与其在制度中的定在、与作为道德规范和法规范之实定化的法律关系。黑格尔认为，这种关系对“一般的法”来说是建构性的。在一定程度上，可以通过分析从法的概念得出，法要实定化，进入定在和实存的外在性之中，并且作为“合法的东西”——法律——而具有效力。哲学的法和实定法彼此是有区别的，但是黑格尔认为自然法和实定法的截然对立是一个“巨大的误解”。他再次以罗马法为例\footnote{《法哲学原理》，第3节，附释，《全集》第7卷，第42页。黑格尔的比较是有问题的，因为不同于黑格尔法哲学教科书的“原理”，制度的原理部分并没有在严格的体系中建立起来。参见F. 舒尔茨（F. Schulz）：《罗马法学史》（Geschichte der Römischen Rechtswissenschaft），魏玛，1961年，第187、194页等处。舒尔茨中肯地对这些制度与亚里士多德的形而上学的松散构造进行了比较。}，拿自然法与实定法之间的相互联系跟《法学阶梯》与《学说汇纂》之间的关系进行比较：《法学阶梯》与《学说汇纂》共属于作为一个实定法体系的《民法大全》。法既由于其在某个国家中有效的形式，也由于其内容的要素，而必然是实定的。“内容的要素”包含了民族的特殊性、一个国家的历史发展阶段以及每一法律体系的组织和适用，这些要素绝不能从任何规范性的法原则中产生出来（第3节）。法必然要走向实定性，由此便产生了对于黑格尔的《法哲学原理》来说乃是奠基性的“自然法”与“国家学”之间的关系，而需要作为纲要来理解的《法哲学原理》一书的副标题也坚持了这种关系。从历史角度来看，这一标题描述的是哲学的两个不同学科。第一个学科（自然法）首先是在近代得到了发展，而第二个学科\footnote{即国家学。——译者}的起源则要追溯到古代世界、希腊的城邦理论。古典政治学体系的优越之处在于，它并不把自然法和国家学分开。柏拉图的《理想国》和《法篇》或者亚里士多德的《政治学》阐述的政治的理性理念，即最佳宪制，把自己理解为对自然法和一种正当地组织起来的（“公民的”）社会（$\pi o \lambda \iota \tau \iota \kappa \acute{\eta}\ \kappa \circ \iota \nu \omega \nu \acute{\iota} \alpha$, societas civilis）的陈述，这种公民社会与国家（$\pi \acute{o} \lambda \iota \varsigma$, civitas, res publica）是同一的。青年黑格尔追随这种从未间断的“国家学”的传统形态，首先在其中看出，对“自然法”进行科学探讨的任务在于，自然法应当在一个政治整体、在“国家”这一基本制度以内建构“伦理自然的法”\footnote{参见《论自然法的科学探讨方式，它在实践科学中的地位及其与实定法学的关系》（1801年），《全集》第1卷，第435页以下。以及本人与之相关的论文：《黑格尔对自然法的批判》，收入《黑格尔法哲学研究》（Studien zu Hegels Rechtsphilosophie），美茵河畔法兰克福，1969年，第42页以下。}。

然而不久之后，黑格尔就不得不认识到，在近代国家的特定历史形态的条件下，应用一种把个体行动埋葬于制度中，由此消除了矛盾的可能性的非辩证的政治概念，只会导致马基雅维利的立场上的政治绝对主义的实践。在与政治绝对主义的分歧中，从霍布斯到康德的自然法形成了与之相对的自然法的个体主义立场。这种自然法的权利在于，要么坚决限制国家对自身秩序井然的社会（“市民社会”）的“非法”干预，要么以一种新的自由概念和法概念革命性地克服已经固化为赤裸裸的权力展示的国家政治。这样，自然法和国家学的对立便引发了向政治启蒙的转向以及向革命的过渡，并且一直伴随着革命的进程。黑格尔的概念辩证法预设了自然法和国家学的对立，并且企图在哲学上和政治上对之予以克服。法哲学把自己理解为“法的哲学”，因为它试图在法理念的辩证法中对这种前国家的自然法与实定的国家法之间的对立进行调和。黑格尔说，法的理念就是自由，要真正把握法的理念，就必须在其概念和概念的定在中（《法哲学》第1节，补充）。也就是说，在各种制度的历史形成过程以及相应的法的现实性的过程中去认识它。

由于法概念与各种制度的定在之间的这种历史的辩证关系，自然法的语言当然就是成问题的了。黑格尔在1817年的海德堡《哲学科学百科全书》中就已经谈道：“迄今为止，自然法这一表达，对哲学法学来说已经成为了一种习惯表达。这一表达具有歧义：法是不是作为一种似乎由直接的自然植入的东西，还是指的是它是由事物的本性，也就是说，由概念来规定自身的。第一种意义是之前习惯所指的意义，以至于同时还虚构出一种自然法在其中应当有效的自然状态。”对黑格尔来说，就法并不是把它自身以及它的所有规定建立在自然之上，而是建立在“自由的人格性——一种为自然规定的对立面的……自我决定”（第415节）之上而言，事物的“概念”就是法。法的现实不是自然；法是“精神的世界，是从精神自身产生出来的”（《法哲学原理》第4节），一种“第二自然”，黑格尔称之为“实现了的自由王国”。在《法哲学原理》概要叙述的那些从最初抽象法的形态开始，一直到家庭、市民社会和国家的各个历史形态为止的制度中，法把自己建立起来。显然，这里处处都有“自然”关系发挥作用。问题在于，这些自然关系是否从自身中产生出了支配那些形态的法律。黑格尔坚决地拒绝这一问题以及对这一问题给予肯定性答复的自然法话语。必然性（比如说，在国家与个体的关系中，归之于“国家”这一制度的必然性）不再意味着，对个体来说，必须生活在国家中，是一条自然法则；相反，国家的伦理必然性以自由的法则为基础。这种自由的法则不是由一成不变的“自然”制定的，而是由——在其辩证法中依赖于各种制度的历史实现的一一概念自身制定的。对黑格尔而言，概念的自我实现乃是人的普遍法权能力的自然法原则，这里所言的人都起来反抗现存的支配和依附关系，并且冲破了迄今为止法的思想、自由都还是隐匿于其下的自然性的假象。

\section*{2}

因此，对黑格尔而言，“哲学法”之有别于国家的实定法，有如它之区别于自然法。在同时代的两种对立的法观念，即康德、费希特的理性法观念与萨维尼的“有机”法观念之间，黑格尔的法哲学处于中间的位置。黑格尔与萨维尼和历史法学派一道拒绝自然法的理性国家。但是，在黑格尔这里，历史解释不仅具有对实定法予以历史论证和辩护的意义，而且具有“概念的发展”的意义（[《法哲学》]第3节）。辩证的概念发展本身就是历史的，即自由在世界历史中的阶梯式进展，而《法哲学原理》的任务就在于，为不同时期的历史——实定法提供“自在自为地有效的辩护”（出处同上）。

这对于理解这部著作及其历史立场有着决定性的意义。黑格尔立足于现代革命及其制宪和立法的基地之上，现代革命的制宪与立法将自然法实定化了。对黑格尔而言，国家的合理性不再是一个公设，而是正在发生的历史现实，“哲学法”承认这一现实。黑格尔提出将“探究理性的东西”与“了解现在的东西和现实的东西”结合起来的哲学概念，反对在他的时代受到偏爱的“关于国家的哲学”，也反对这种国家哲学炮制出新的和特殊的理论的目标，

\begin{quote}
“似乎世界上从未有过国家和国家制度，现在也还没有，但是现在——这个现在是永远继续下去的——似乎应该从头开始”\footnote{《法哲学原理》，序言，第23页。（中译文引自《法哲学原理》，第4页。——译者）}。
\end{quote}

黑格尔承担了哲学的传统任务“理解存在的东西”，也就是对存在的思辨认识，以对哲学提出一种似乎将其颠倒了过来的解释：在时间发生的视域中思考存在。因为哲学作为对存在的东西\footnote{指哲学。——译者}的认识，同时又是“在思想中得到把握的它自己的时代”；哲学无法跳出它的世界，就像个人作为“时代之子”无法跳出罗陀斯岛一样。

黑格尔将“法的理念”规定为概念和定在的统一，就是基于时间和存在的这种交织。在解决将自由认识为这种理念这个问题的过程中，《法哲学原理》避免了自然法和法实证主义的双重片面性，它们要么仅仅执着于“自然”法的抽象概念，要么执着于实定法的定在。理念“为了得到真正的理解，必须在法的概念及其定在中来认识法”（第1节，“补充”）。其中包含了理念与历史的联系。自由只有在这样的前提下，即它已经在历史中进入定在并且内在于实定法之中了，才能作为法的理念得到把握。这发生在现代世界的国家中。现代世界的国家在自身中接受了法国大革命的自由原则，并且具有保存个体为自由的、具有法权能力的人格的任务。黑格尔把现代国家定义为“逻辑学”思辨意义上的“现实”，成为直接之物的内在本质与外在实存的统一：现代国家是“伦理理念的现实”（第257节）、“实体性意志的现实”（第258节）以及“具体自由的现实”（第260节）。这样，不同于从卢梭到费希特为止的观念论自然法，对黑格尔而言，法理念并不是思维的单纯可能性，而是一种经过历史的中介的现实，这种现实要求哲学放弃其抽象的自然法立场，法的理念要求哲学把当前的国家作为“其自身是一种理性的东西来理解和叙述”\footnote{《法哲学原理》，序言，第34页。（中译文引自《法哲学原理》，第12页。——译者）}。

黑格尔在《法哲学原理》序言中的格言：“凡是合乎理性的东西都是现实的；凡是现实的东西都是合乎理性的”，所指的便是这种现代世界的国家，而不是1821年的普鲁士国家。要认识这句格言的“进步”意图，我们不需要用第一句来反对第二句：因为这两句说的是一样的：在现代国家中，法的概念已经达到了定在，理性的东西——自由的思想——是现实的了，并且现实的东西——现代世界的国家——是理性的了。这并没有排除部分的不理性。黑格尔在最后一次法哲学讲演（去世前几天）对这句格言的一个解释中，简洁明了地指出：“现实的东西就是理性的。但是并不是所有实存的东西都是现实的，坏的东西是一个自身败坏的和虚无的东西。”\footnote{《法哲学讲演录》第2卷，伊尔廷出版，斯图加特，1974年，第923页。}

因而，黑格尔在分析法国大革命的自由原则时，提出了哲学的实现问题。那在法国大革命中起而反对已经成为无精神、无理性的封建制度的现实——“不法的旧制度”——并将其推翻的法的概念起源于哲学思想。而这也是这一事件无与伦比和史无前例之处。对于黑格尔而言，这一事件具有世界历史的意义：

\begin{quote}
“自从日立中天、众星绕行以来，还从未见过这样的事情，人立于头脑——也就是说思想——之上，并据此建造现实性。阿那克萨哥拉最先说出，努斯统治着世界；而只是到了现在，人才进而认识到，思想应当统治精神的现实性。”\footnote{《世界历史哲学讲演录》第4卷，G. 拉松出版，莱比锡，1920年，第926页。}
\end{quote}

但是哲学的实现并不意味着哲学的扬弃，并不是思想消解于一个历史上改变了的现实性之中，而是相反：要理解这种改变的根据，就需要哲学。因而哲学就成为了时代的理论，这是在直到黑格尔为止，哲学在其历史中从未有过之事；哲学取得了这样的任务——为了拯救法国大革命的自由原则：(1)反对其自身落入暴政之中；(2)反对复辟。复辟与当下的原则直接对立，因而哲学否定了复辟想要重新建立的历史连续性\footnote{参见里特尔（J. Ritter）：《黑格尔与法国大革命》（Hegel und die Französische Revolution），美茵河畔法兰克福，1965年，第24页以下。}。在这个意义上，黑格尔哲学是革命的哲学，而不是复辟的哲学或者普鲁士国家的哲学（海姆）。黑格尔哲学不是以改变的意志，而是以从法的现实性的事实中获得其原则的认识意志，面对时代及其“现实”。黑格尔在暗中提到培根的一句话时不无道理地说，半吊子哲学（尤其是哈勒、施莱格尔、德·迈斯特、博纳尔德的浪漫主义复辟理论属于其列）离开了国家，真正的哲学则导向国家\footnote{参见《法哲学原理》，序言，第36页。（中译文参见《法哲学原理》，第13页。——译者）}。这种“真正的哲学”的缺陷或许在于，它只是想要认识，并且通过认识，不是想要改变现实，而是想要“和解”。这种哲学的优越之处首先在于，它在概念——内在于完成了的历史之中的概念——的媒介中，对一种改变了的现实做出解释。

\section*{3}

因此，《法哲学原理》是那种青年黑格尔在《精神现象学》中提升为他的哲学规划的“被概念把握的历史”的一个教育剧本。对黑格尔而言，概念就是在其自由中的人的意志，自由构成了意志的“实体和规定”（第4节）。法在概念的自我展开中的辩证论证既涉及古典政治学和现代自然法的诸范畴，也涉及历史—精神的“现实”的形成过程，通过这一过程，这些范畴获得了新的内容。这一点从（“法”的诸制度的）体系纲要以及体系的三分法（一、抽象法；二、道德；三、伦理）可以看出来。抽象法是18世纪的自然法(jus naturale absolutum et hypotheticum)，探讨的主要是私法上的内容（财产、契约）；它是从一切支配和同业公会的束缚中解放出来了的抽象个体的法。黑格尔按照自然法的模式，把它放在《法哲学原理》的最前面，并且显然修正了他最初的立场，在不依赖于“实体性伦理”（家庭和国家）的前提下，对它进行阐述。对黑格尔而言，不是“具有一定身份而被考察的人”，不是作为一个不同于奴隶、农奴或“亲属”的市民，而是具有法权能力的“人格本身”（第40节）才是法和财产概念的基础。道德包含了经院哲学在伦理学的名义下进行探讨的那些概念和内容：古代意义上的德行论和义务论，但是黑格尔却将经由康德的主体性转向所形成的反思道德的原理纳入其中。

与经院传统相一致，黑格尔把抽象法（自然法）理解为“法的意志”的学说，把道德（伦理学）理解为“道德意志”的学说。在这两种学说中，随着现代革命——既是哲学革命也是政治革命——而立足于自身之上的人之为人，一方面在与外部的关系中，在与外部自然世界和周围人类世界的关系中，以财产（第41-71节）、契约（第72-80节）和不法（第84-86节）的形式，另一方面则是在与内在性的关系中，以故意和责任（第105-118节）、意图和福利（第119-128节）、善和良心（第129-141节）的方式，实现了自身，并将自己表达出来。尽管黑格尔承认这些范畴，并且把这些范畴吸收到了他的演绎中，但是这些范畴对黑格尔而言仍是抽象物，在《法哲学原理》体系中只具有有限的效力。被标明为“抽象法”的凝结为私法的自然法，乃是具有法权能力和财产能力的一般人格的法，借由康德-费希特哲学而成为主观性和反思性的伦理学也同样是与自身相关的主体性的道德。

如果说抽象法和道德与近代自然法和道德学(Moral)的范畴之间有一种显而易见的平行关系，那么如果把这种比较运用到《法哲学原理》的第三部分，就会产生不少困难。这些困难不仅仅在于家庭、市民社会和国家的划分，而且也在于伦理这个名称本身。因为在17和18世纪，自然法也划分为个体法和社会法，前者以财产和契约为内容，后者以家（家庭）和市民社会(societas civilis)或者国家(civitas)为内容。黑格尔的伦理理论仍然意味着把一个可以追溯到古典政治学的范畴重新引入《法哲学原理》之中。对于现代自然法而言，伦常(Sitte)并不具有一种独立形式的实践生活的效力，而是一般地从属于法概念的合理性之下\footnote{参见博比奥（N. Bobbio）：《黑格尔与自然法学说》（Hegel und die Naturrechtslehre），1966年，重刊于M. 里德尔（出版）：《黑格尔法哲学资料集》（Materialien zu Hegels Rechtsphilosophie）第2卷，美茵河畔法兰克福，1975年，第81页以下。}。这也适用于康德的法权学说。众所周知，法权学说是康德所称的《道德形而上学》的一部分。与之相反，黑格尔把伦理范畴与伦常的传统立场重新结合起来，在伦理范畴中，法国18世纪作家（孟德斯鸠）对“mœurs”（风俗）的发现，对柏拉图-亚里士多德城邦理论的重构以及早期浪漫派的“民族精神”的“有机”统一的理论相互交织在一起。黑格尔的“伦理”指的是个体与“伦理力量”（第145节）以及每一特定的、历史的民族和国家的“必然关系”（第148节）的统一。因此，伦理也就代表了道德学(Moral)与政治学(Politik)的关系，对于国家与“市民社会”统一(civitas sive societas civilis)的传统理论来说，这种关系是根本性的。黑格尔在这里提到“客观的”而非“道德的”（亦即“包括在道德主观性的空洞原则中的”）“伦理的义务论”与“第三篇中对伦理必然性领域”、国家内部的伦理关系和制度的“系统阐述”之间的同一性，并不是没有道理的（第148节）。

然而，制度论与伦常传统的结合仅仅触及黑格尔伦理理论的一个方面。另一个方面是，这个在形式和内容上超出近代自然法的法哲学范畴记录了从18世纪转向19世纪时在政治社会世界建构中已经完成的根本变化。毫无疑问，这里，对于各种制度的体系性重构遵循了《逻辑学》的概念辩证法。在《法哲学原理》中，黑格尔处处以《逻辑学》的概念辩证法为前提。尽管“概念”中包含的个体性、特殊性和整体性的三个环节已经为1817年的《哲学全书》中伦理的划分奠定了基础，但此时“民族”这一基本制度还是一种静态的统一，尚未按照特殊和普遍的环节进行区分，而是作为整体代表个别性范畴\footnote{《哲学全书》，[第三部精神哲学，]二、客观精神，3. 伦理：(1) 单个民族；(2) 外部国家法[在和平与战争状态中相互对立的（众多）“特殊”国家意义上的特殊性]；(3) 普遍的世界历史（即从特殊的“民族精神”的辩证法中产生出来的“普遍”精神意义上的普遍性）。参见《哲学全书》，第430-452节。}。黑格尔在《法哲学原理》中放弃了这种立场，而让“个别民族”的实体性伦理再度服从概念的辩证法：它分化为家庭（个别性）、市民社会（特殊性）和国家（普遍性）。

《法哲学原理》第三部分的这种划分涉及一个悠久的、令人崇敬的传统。黑格尔接受了这一传统，以便借助辩证法的手段消解这一传统，并在他的思想中克服之。黑格尔辩证地消解传统范畴的最重要理由之一就是要引入一个全新的社会概念，他由此开创了国家哲学历史的新纪元。直到今天，黑格尔借由这一概念让时代及其当下所意识到的，绝不亚于现代革命的历史成就：一个去政治化的、建立在个体的自由和平等之上的“市民社会”的产生，随着从英国开始的工业革命，市民社会的重心也从政治组织形式转移到了经济上面。自由主义自然法的社会思想并没有或者说仅仅是不充分地关注到国民经济学的语言，而在黑格尔这里，国民经济学则构成了一种从根本上发生了改变的概念演绎的参照系。按照自然法的观点，一切“社会”都是人格性地组织起来的，都是由许多个人构成的团体，他们通过理性的言论和由理性规定的行动形成一个对所有人都有法的约束力的共同意志，也就是一个“法权人格”的意志。这里，言和行都具有公共的意义，并在将每一个“社会”（而不仅仅是国家）建构为法律制度(Rechtsfigur)的那种公共领域的媒介中实现出来。相反，对黑格尔来说，“社会”按照定义是由通过需要和劳动彼此联系在一起的私人构成的。劳动是一种特殊的行动方式，需要是作为“私人”之人的自然基础。黑格尔在这一领域的哲学成就首先在于，将其“私人的”参照系理解为经过了公共中介的，并且将其自然基础理解为社会的常量。“市民社会”理论的定位系统并不是契约，即理性的、言行出众的个人之间的联合，而是“需要的体系”，也就是从需要、劳动以及满足需要的手段中产生出来的且在其活动中不断再生的“私人”（第187节）之间的关系网络。

黑格尔“辩证地”理解的私人目的与公共目的之间的交互关系，作为一种超越于“个人”行动之上的“社会联系”的基础，已经在英格兰-苏格兰道德哲学（休谟、斯密、弗格森）和法国百科全书派的利益理论中反映出来。在这里，通过将私人目的、对个人生命和财产的保护宣布为公共事务，契约模式就在其形式法功能之外获得了一种内容。功利和利益理论支配了欧洲晚期启蒙运动的思想，青年黑格尔研读了欧洲晚期启蒙运动的那些文本。它们宣扬的是市民私人彼此之间不受等级和地产关系的限制而进行交往的原则，这种交往通过“个人利益”（狄德罗、爱尔维修、霍尔巴赫）或者“自利”（富兰克林、边沁）规范了所有的社会关系。然而，它们和谐的社会图景现在被解释为冲突模式。黑格尔对那个不再与其他概念（国家、市民社会）同义的行为体系的逻辑-历史问题的态度是模糊的。一方面，“社会”是一个法的构造物(Gebilde)，是理性意志（“自由意志”）的定在的法；另一方面，意志只在私法领域才将自己概括为契约——社会，作为在家庭、市民社会和国家中实现其自身的自然-历史制度，并不是建立在契约的法律义务上的\footnote{《法哲学原理》，第181节，《全集》第7卷，第261页以下。黑格尔认为契约模式不足以从概念上把握家庭（第75节，第163节）和国家（第258节），而予以拒斥。}。

尽管“市民社会”也表现为法的构造物，但对黑格尔来说，在从理性意志（“自由意志”）的理念出发对其进行辩证的论证和辩护时，真正重要的是要表明这个制度落在了传统自然法的“社会”的起源以外，因为一方面它在体系上和历史上远远超出“家庭”这一制度之外，另一方面又没有达到“国家”这一制度。在[市民社会]这一阶段，被黑格尔普遍化为“法的理念”并整合到（客观）精神的辩证法之中的理性意志，将自己“特殊化”为许多异质的个人，“因为这些人本身在他们的意识里所有的和成为其目的的，不是绝对的统一，而是他们自己的特殊性和他们的自为存在，——这就是原子论的体系”\footnote{《哲学全书》，第三部分，第523节，《全集》第8卷，第401页。（中译文参见《哲学科学百科全书III：精神哲学》，第291页。——译者）}。

在这种制度的辩证法中构成的东西，是一种区别于国家、在需要和劳动的基础上成长起来的并且在其活动中不断再生的“私人”之间的关系网络——现代词义上的“社会”。它作为“市民社会”进入法哲学及其政治理论的中心，在这种地位中熔化了古典政治学和现代自然法的那些传统范畴。这首先表现在，社会并没有局限于划归给它的地方，而是扩展到了其他制度领域。因为，在黑格尔那里，市民社会一方面是抽象法的“定在”（第209节以下），另一方面又是主观道德的活动领域（对贫困领域来说）（第242节）、家庭的实体性基地和力量（第238节），最后还是国家的差异（第242节）。在伦理的建构中，市民社会以在历史上得到规定的方式成为了“中心”。在这里，市民社会表现为介于“家庭”与“国家”这两种制度之间的连接环节，这在古典政治学和现代自然法传统中是付之阙如的。市民社会这一“中心”不仅是在平衡与和解各种对立意义上的中介，而且更是形成了一个紧张的领域，这个领域把所有的制度都纳入其影响范围，并且改变了它们的传统结构。

这始于“家庭”这一制度，即“市民社会”的基本元素，按照自然法和国家学的传统理解，家庭乃是其自身构成为政治国家的中介。黑格尔之所以剥夺家庭的这一中介功能，不仅仅是因为市民社会和国家分离开来，而且也因为家庭自身按照其“概念”不能是国家团体的一个环节。家庭与国家的关系仅限于家庭是国家的“材料”（第262节）。也就是说，既在生命的自然繁衍过程中准备好作为“群众”的个人，也在教育中准备好直接的伦理情感和意识（第173-175节）。总而言之，在黑格尔之前，“家”之制度作为“整个家庭”，本身就是国家整体的一分子，现在，这个分子[不再是家庭，而是]在家庭的私人领域中成长起来并受到教育的个体。与此相应的是制度的构建过程中的一种与历史相平行的辩证法。一方面，与国家的关系要从家庭过渡到市民社会——它是黑格尔新颖地嵌入伦理体系之中的制度领域；另一方面，市民社会和家庭也必须将它们自己“设定”到一种制度上改变了的关系中。对黑格尔而言，家庭并不是一个由具有不同身份的成员所构成的“社会”，亦非由“更小的社会”（丈夫-妻子、父母-子女，主人-奴隶）组成，而是表现了一种“人格”，它在“一种所有权”中，而不是在家庭的法-经济上的统一中，具有其外在实在性（第169节）。在黑格尔这里，家庭概念的原初的“家政”特征已经消失殆尽，被一种现代的情感特征取而代之。这种情感特征完全建立在“私人”关系和情感的基础上，在18世纪末开始形成。因此，第158节这样定义家庭的概念：“作为精神的直接实体性的家庭，以爱为其规定，而爱是精神对自身统一的感觉。因此，在家庭中，人们的情绪就是意识到自己是在这种统一中，即在自在自为地存在的实质中的个体性，从而使自己在其中不是一个独立的人，而成为一个成员。”\footnote{中译文引自《法哲学原理》，第175页。——译者}黑格尔表述这一概念的时候完全认识到了《法哲学原理》体系中的市民社会的基本地位所引起的传统家政学形态的改变。因为，与迄今为止的政治学和自然法中的市民社会(societas civilis)理论——按照这些理论，市民社会以“家庭”社会为基础——相比，家庭在黑格尔所讨论的“市民社会”中是一个从属性的东西，仅仅是奠定基础而已：家庭不再具有如此“全面的影响力”。在超出了“整个家庭”界限的现代社会的经济再生产过程的条件下，个体已经从家庭中被“揪了出来”，并且如黑格尔所言，成为“市民社会的子女”；市民社会接过了经济的功能，“又用它自己的土地来代替外部无机自然界和个人赖以生活的家长土地，甚至使整个家庭的存在都依从它，而听偶然性的支配”（第238节）\footnote{中译文引自《法哲学原理》，第241页。——译者}。

正如黑格尔的市民社会概念一方面涉及改变了的家庭结构，另一方面它又涉及同样变化了的国家的地位。当我们把黑格尔伦理理论的第三章与传统政治学理论进行比较的时候，就可以确定，由于其种种前提，它不再是对“市民社会”及其政治制度（政府）的陈述，而是与其形成鲜明对照的“国家科学”。在黑格尔这里，传统通过市民社会(societas civilis)的关系概念融为一体的国家和市民社会第一次“设定”为一种关系，而且设定为分离、“差别”的关系（第182节，补充）。因而黑格尔拒绝“宪制的古老区分”，因为他认识到，在“差别”的前提下，奠定这种区分基础的社会和政治制度的“实体性统一”并不是“理性的”，或者按照黑格尔《法哲学原理》的前提所指的同样的东西，不再是“现实的”。它必然沦落为一种历史性的、只合乎“古代世界”立场的现象：“古代把国家制度区分为君主制、贵族制和民主制，这种区分是以尚未分割的实体性的统一为其基础的，这种统一还没有达到它的内部划分（一个在自身中发展了的机体），从而也还没有达到深度和具体的合理性。因此，从那种古代的观点看来，这种区分是真实的和正确的……”（第273节）\footnote{中译文引同上书，第287页。——译者}。

然而，[黑格尔]法哲学与传统的断裂不仅仅表现在，由于国家和市民社会的差别，以前的政治学概念和内容已经变成非理性的，成为了历史的逝去之物；通过为黑格尔的伦理理论奠定基础的国家与历史间的交互关系，这个决裂更是无比清晰地呈现出来。黑格尔的国家概念，往后看，以古代市民社会的“实体性统一”的瓦解和[新]出现的与市民社会的差异为前提；往前看，则与更广阔的世界历史领域相连。黑格尔使市民社会位居国家之下，与他使历史位居[国家]之上，两者形成一种互补的关系。如果我们看一下历史在迄今为止的政治学和自然法理论中的价值，就可以说，两者的共同做法是，都把法的问题或者“好生活”的问题化解为“最佳宪制”问题，也就是说建“国”(civitas)问题。古代人和现代人都认为，“最佳宪制”是与市民社会(societas civilis)、正当地组织起来的并且使得德行的施展得以可能的生活的建立同时发生的。为此，只需要回顾“先前的”科学，如形而上学、人类学、心理学，以及作为德行论和义务论的伦理即可。但是，所有这些科学都只不过构成了开端，作为市民社会的国家的学说基础。在18世纪，随着诸如国民经济学和历史哲学这些新学科的产生，前述科学在实践哲学上的紧密结合也就随之瓦解了。即便是最早在法学和国家学问题与历史哲学问题之间建立起联系的康德，仍然以传统自然法的方式把这个主题阐述为“达成一个普遍地管理法权的公民社会”\footnote{《在世界公民底观点下的普遍历史之理念》，[《全集》]第8卷，科学院版，第22页。（中译文参见〔德〕康德：《康德历史哲学论文集》，李明辉译，联经出版事业公司2002年版，第11页。——译者）}。而黑格尔则不同，在他看来，历史哲学把自身从自然法与“市民社会”理论的联系中解脱了出来。

国家的普遍理念在其完成的顶端下降为“众多国家”的特殊性，在历史的游戏中，这些国家彼此对立。国家理念是“作为类和对抗个别国家的绝对权力”的理念，是“在世界历史的过程中给自己以现实性的精神”（第259节）。在黑格尔这里，国家不再是传统的政治状态概念，不是古典政治学和现代自然法的无运动的、非历史的自然或契约模式，而是历史——作为实践哲学的新法庭——屹立于其上。在历史中，“伦理性的整体本身和国家的独立性”——鉴于家庭的伦理性整体所遭受到的同样命运，我们必须为黑格尔补充道：再次——“都被委之于偶然性”（第340节）。黑格尔在《法哲学原理》的末尾所引入的历史维度，是自然法学家在其体系的开端处设定的那个自然状态的理念的现实性。在自然法学家那里，运动乃是从自然状态走向“societas civilis”意义上的市民社会，并在这里进入静止状态；在黑格尔这里，则是当国家不再关系到市民社会，而是关系到其他国家的时候，这种运动便开始了。这种自然状态是一种现实的、绝非虚构的状态，是历史的运动，《法哲学原理》把这种历史运动接受到自身之中，借此再一次从抽象的法和社会的自然理论中解放出来。法概念的自我展开从“普遍精神”的定在要素开始，并结束于它。在世界历史中，普遍精神的定在要素就是“精神的现实性”，在其中人和物的“自然”并不是一种永恒的规律，而是黑格尔思为自由的概念，它“在其内在性和外在性的全部范围内”实现它自己。精神的运动与法哲学结束、历史哲学开始同时发生，在这一运动的普遍性之前，“只是作为观念性的东西……”的家庭、市民社会和国家，“并且在这个要素中，精神的运动就在于把这一事实展示出来”（第341-342节）\footnote{中译文参见：《法哲学原理》，第351页。除了“观念性的东西”外，引文后面的强调为里德尔所加。此外，引文来自第341节，而不是来自第341-342节。——译者}。

\section*{4}

如果对黑格尔《法哲学原理》的结构予以总结性地考察，那么我们可以说，那些具有历史意蕴的概念形式为逻辑上的各个建构原则奠定了基础。那些伦理事物的制度，特别是“社会”这个概念，还有“世界历史”的终结概念，都是逻辑的—历史的范例。这两个概念同时改变了古典政治学和现代自然法数世纪以来一以贯之地予以阐述的家庭和国家的体制。就像家庭处于作为成员而从属于它的个人与市民社会之间一样，国家也运行于市民社会和历史之间。在黑格尔这里，正是传统政治制度结构上的这种缝隙，推动了辩证地交织在一起的法概念的发展历程及其体系的精巧的结构形式。由此出发，《法哲学原理》尤其获得了其所特有的丰富的历史内容，正是这一点把它从原则上与近代自然法著作区分开来。

然而，《法哲学原理》作为制度辩证法的体系引入现代世界政治理论之中的这些新元素，也成为了黑格尔身后那些理论的出发点，这些理论动摇了他的体系，突破了他的那些中介。对《法哲学原理》的批判，正好就是挑战黑格尔本人在法的现实性的哲学重构中对逻辑的东西与历史的东西所提出的中介：(1)国家和历史的关系和(2)市民社会和国家的关系。对于黑格尔《法哲学原理》（1833年）的最早出版人爱德华·甘斯(Eduard Gans)来说，该书最重要的价值在于，黑格尔在书末为我们描绘的“壮丽景象”：他让国家“卷入历史的世界海洋”。甘斯不无道理地指出，黑格尔对于向历史的过渡所做的简要勾勒“不过预示了”这一领域中的“更为重要的关切”\footnote{《法哲学原理》，前言，《全集》第8卷（1833年），第VII—IX页。}。由此，如同他自己（以及后来的黑格尔主义者费迪南德·拉萨尔）已经用他的《世界历史发展中的继承权》（1824-1835年）向历史主义思维方式致敬一样，甘斯预见到了支配了19世纪的历史主义对哲学的未来批判。历史主义的本质在于，它不再把国家理解为自身的现实，而是理解为历史的产物。基于这种观点，阿诺德·卢格（1842年）对黑格尔的法哲学进行了批判。在卢格看来，《法哲学原理》并不关心时代和历史的进步运动，因为它在前述意义上定义国家。然而，国家实际上并不应当是普遍的现实，而是如同卢格从历史出发对其进行批判的1840年的普鲁士国家一样，是一个个别的—历史的实存。尽管《法哲学原理》中的国家令人想到“当下的国家”，是的，黑格尔甚至指出这个国家的名称，但他并没有让它从“历史过程”中产生出来，这也是为什么黑格尔对政治的生活意识和时代意识的发展毫无影响的原因所在\footnote{卢格：《黑格尔法哲学与我们时代的政治》（Die Hegelsche Rechtsphilosophie und die Politik unserer Zeit），《德意志年鉴》1842年卷，第762页以下。}。

批判的第二个出发点就是市民社会和国家的关系，这是青年马克思（1841或1843年）评论《法哲学原理》第261-313节的主题。马克思颠倒了黑格尔对市民社会和国家的关系的思辨解释：国家是市民社会的现象，而黑格尔贬为“伦理事物的现象世界”（第181节）的市民社会，那种远为坚实的、同时完全非思辨地理解的政治经济学的现实，对马克思来说，才是理解国家和历史的关键。马克思由此取消了黑格尔法哲学的中介。在这里，国家也就从一种现实贬低为一种产物，也就是经济的产物。政治经济学批判的问题不再是市民社会的“私人”原则与政治国家的中介，而是对这个市民社会及其私人—个人主义原则本身的批判。马克思抛弃《法哲学原理》，转向英国人和法国人的政治经济学。这并没有妨碍他对那种在黑格尔市民社会理论中开始浮出水面（第243-245节）的贫富辩证法予以批判。通过将[贫富]这两个极端及其基础（黑格尔发现的市民社会领域）独立出来，马克思将辩证法驱至一场[新的]运动，它超越了《法哲学原理》从个别的个体经过家庭和市民社会再到国家的概念运动。按照将马克思与圣西门、孔德、约翰·斯图亚特·密尔和斯宾塞联系起来的社会学的思维方式，国家是社会运动的产物，而不是黑格尔赋予它的那种伦理“理念”的现实性。

就此而言，《法哲学原理》实际上处于一个终结点，但它并没有那种产生于工业社会早期阶段的“国家式微”的思想。它只是通过把国家和市民社会分离开来，为这一思想的出现做好了准备，但却未能设想到这一思想。黑格尔哲学概念还包含了限制的智慧，因此他不可能通过其哲学概念做到这一点。因为哲学既不想教导世界应当怎样，也不能预言世界将会怎样。在黑格尔看来，哲学作为世界的思想，只有在历史形成过程完成和“结束”的时候，才会显现在时间之中。因此，法哲学执着于一种“旧”的生活形态，这种形态可以认识，但是不能使之变得年轻或者进行改变。认识的任务就在于，不要让哲学超出它自己的时代界限。黑格尔在《法哲学原理》序言的最后说道，密涅瓦的猫头鹰只有在黄昏来临之后才起飞。但这并不意味着，哲学对我们而言就属于过去，其本身已经成为“历史的”了。时下，东方和西方之所以同样都对《法哲学原理》进行深入的研究，是因为工业社会的现实在转入其对立面之前，绝对不会兑现圣西门和马克思的那些超越了它的发展的思想。在现代社会的形成过程中，国家变成了一种神话般的存在，它在自身中将权力和智慧统一了起来，并且剥夺了具有矛盾可能性的个体的自由，而在黑格尔看来，自由乃是思想的源泉。意识形态和伟大预言的时代并没有阻挡住其所反对的政治统治的神话，反而是促进了它。利维坦的回归是向哲学提出的一项挑战，哲学必须接受这一挑战，以便能够重新找回它的经典志业：通过理性的力量摧毁神话。黑格尔的《法哲学原理》就是这种力量的见证。绝非偶然的是，人们已经再一次从黑格尔开始，在一种“自身”从政治统治中解放出来的基地上对自由、法和国家的关系进行讨论。今时今日，澄清这种关系要比以往更加紧迫。